\documentclass[11pt]{amsart}
\usepackage[margin=1.1in]{geometry}
\usepackage{amsmath,amssymb,amsthm}
\usepackage{booktabs}
\usepackage[colorlinks=true,linkcolor=blue,citecolor=blue,urlcolor=blue]{hyperref}

\newtheorem{theorem}{Theorem}
\newtheorem{proposition}[theorem]{Proposition}
\newtheorem{lemma}[theorem]{Lemma}
\theoremstyle{definition}
\newtheorem{definition}[theorem]{Definition}
\newtheorem{hypothesis}[theorem]{Hypothesis}
\theoremstyle{remark}
\newtheorem{remark}[theorem]{Remark}

\newcommand{\lbl}[1]{\textup{\textsc{#1}}}
\newcommand{\DHL}{\mathrm{DHL}}
\usepackage[htt]{hyphenat}
\newcommand{\tl}[1]{\texttt{#1}}
\emergencystretch=4em
\raggedbottom

\title[The finite part of Maynard's reduction in core Lean]{The finite part of Maynard's reduction from a variational ratio to bounded gaps, checked in core Lean}
\author{Jerry Kim}
\thanks{Prepared with the assistance of Claude (Anthropic). The Lean sources are in \url{https://github.com/jerryckim16/twin-primes} (directory \tl{lean/}); the project's main paper is at \url{https://jerryckim16.github.io/twin-primes/}.}
\date{Version 1.0, 26 November 2026}

\begin{document}

\begin{abstract}
Maynard's method proves bounded gaps between primes in two kinds of steps:
analytic asymptotics for two weighted sums $S_1$, $S_2$, and finite,
combinatorial steps around them. This note records which finite steps we
have checked in Lean~4 without Mathlib, and lists precisely the analytic
statements that stay hypotheses. Checked (\lbl{Lean-checked}): the pigeonhole
step from a sieve inequality to $\DHL[k,2]$; the passage from a ratio above
$4$ to a level below $1/2$; an abstract limit step from two-sided asymptotics to
the sieve inequality; for $k$-dimensional weights, the exact expansion of
$S_1$, the Chinese-remainder count with its error, Maynard's change of
variables $\lambda\leftrightarrow y$ for every $R$ and $k$, the
diagonalisation of the main term, and the exact contribution of pairs whose
moduli are not pairwise coprime. Not formalised: the asymptotics of $S_1$ and
$S_2$ (which rest on the Bombieri--Vinogradov theorem), the size of the error
terms, the finite expansion of $S_2$, and the variational inequality
$M_k>4$. A Mathlib formalisation of $H_1\le246$ by another group exists; this
note is a small independent check of the combinatorial skeleton, not a
formalisation of bounded gaps. Nothing here bears on the twin prime
conjecture, which remains open.
\end{abstract}

\maketitle
\setcounter{tocdepth}{1}
\tableofcontents

\section{Introduction}

\subsection*{The reduction}
Let $\mathcal H=\{h_1<\dots<h_k\}$ be admissible. Maynard \cite{May15} takes
weights $w_n=\big(\sum_{d_i\mid n+h_i\,\forall i}\lambda_{d_1,\dots,d_k}\big)^2$ on
$n\equiv v_0\pmod W$ and compares
\[
S_1=\sum_{\substack{N\le n<2N\\ n\equiv v_0\ (W)}} w_n,\qquad
S_2=\sum_{\substack{N\le n<2N\\ n\equiv v_0\ (W)}} w_n\,\#\{i:\ n+h_i\ \text{prime}\}
\]
(\lbl{Literature}, \cite[(4.2)--(4.3)]{May15}). If $S_2>S_1$ for all large $N$, some
$n$ has two primes among $n+\mathcal H$. His Proposition~4.1 evaluates $S_1$ and
$S_2$ asymptotically in terms of a test function $F$, and Proposition~4.2
turns $\theta M_k/2>1$ into $\DHL[k,2]$ \cite[pp.~5--6]{May15}.

\subsection*{What this note does}
We separate the finite steps from the analytic ones and check the finite
steps in Lean~4.35.0, core library only (no Mathlib). Every theorem named
below compiles with no \tl{sorry}, no \tl{native\_decide} and no new axioms;
\tl{lake env lean Audit.lean} reports only \tl{propext}, \tl{Quot.sound} and
\tl{Classical.choice} (Section~\ref{sec:audit}). Section~\ref{sec:hyp} lists
the analytic statements that are \emph{not} formalised. In core Lean they
appear only as hypotheses of theorems; none is assumed as an axiom.

\subsection*{Relation to other work}
A large Lean~4 formalisation of $H_1\le246$, built on Mathlib and
PrimeNumberTheoremAnd, has been announced \cite{Axiom26}. Its site says it
formalises ``the majority of Maynard's paper'' and the part of Polymath8b
\cite{Pol14b} used for $246$ (\lbl{Literature}, read 26 November 2026; we have
not inspected its code or its remaining hypotheses). Our development is
independent of it and much smaller. Its value is only as a second,
self-contained check of the combinatorial skeleton; the statements below are
standard and possibly folklore in this form.

\section{The skeleton: from a sieve inequality to bounded gaps}\label{sec:skeleton}

All objects are natural numbers or finite lists. $\tl{IsPrime}\ n$ means
$n\ge2$ with only the trivial divisors; $\tl{Admissible}\ H$ means $H$ misses a
residue class modulo every prime; $\tl{DHL}\ k$ says that for every admissible
increasing $k$-tuple $H$ and every $N$ there are $n\ge N$ and $a<b$ in $H$ with
$n+a$, $n+b$ prime (file \tl{Basic.lean}).

\begin{definition}[\tl{SieveInequality k}, file \tl{Reduction.lean}]
For every admissible increasing $k$-tuple $H$ and every $N$ there are a finite
list $ns\subset[N,\infty)$ and weights $w:\mathbb N\to\mathbb N$ with
$\sum_{n\in ns}w(n)<\sum_{n\in ns}w(n)\,\#\{h\in H:\ n+h\ \text{prime}\}$.
\end{definition}

\begin{theorem}[\lbl{Lean-checked}]\label{thm:skeleton}
\begin{enumerate}
\item \tl{dhl\_of\_sieveInequality}: $\tl{SieveInequality}\ k\Rightarrow\tl{DHL}\ k$.
\item \tl{sieveInequality\_of\_ratio}: if for all $H$, $N$ there are weights with
$S_1>0$ and $pS_1\le qS_2$ for fixed $q<p$, then $\tl{SieveInequality}\ k$.
\item \tl{theta\_of\_ratio\_gt\_four}: if $p>4q$ there are $a,b$ with
$2a<b$ and $2bq<ap$, i.e.\ a level $\theta=a/b<1/2$ with $\theta(p/q)/2>1$
(take $a=p+4q$, $b=4p$).
\item \tl{eventually\_lt\_of\_asymptotics}, \tl{sieveInequality\_of\_asymptotics}:
if $a c'<c a'$, $B(N)>0$, and for every $m>0$ eventually
$m a' S_1(N)\le(m+1)aB(N)$ and $m cB(N)\le(m+1)c'S_2(N)$, then eventually
$S_1(N)<S_2(N)$ (one may take $m=2ac'+2$); applied to every tuple this gives
$\tl{SieveInequality}\ k$.
\item \tl{boundedGaps\_246\_of\_asymptotics\_50}: the hypothesis of (4) for
$k=50$ implies $\tl{BoundedGaps}\ 246$, through an admissible $50$-tuple of
diameter $246$ checked by kernel reduction (\tl{Tuples.lean}).
\end{enumerate}
\end{theorem}

\begin{remark}
In (4) the constants $a/a'$ and $c/c'$ play the parts of $I_k(F)$ and
$(\theta/2)J_k(F)$ in Maynard's normalisation. Only the upper bound for $S_1$
and the lower bound for $S_2$ are used. Everything is in $\mathbb N$, so the
real constants of \cite{May15} must first be replaced by rational bounds on
the correct side; this costs nothing in the argument but is not formalised.
\end{remark}

\section{The finite part of the evaluation of \texorpdfstring{$S_1$}{S1}}\label{sec:s1}

Tuples are lists; $d\mid n+h$ means $d_i\mid n+h_i$ for all $i$; the support is a
finite list $D_k$ of tuples and $\lambda:D_k\to\mathbb Z$. Lean's $S_1$ runs over
$n<N$; Maynard's range $N\le n<2N$ is the difference of two such sums.

\begin{theorem}[\lbl{Lean-checked}; files \tl{MaynardS1.lean}, \tl{MaynardHpw.lean}]\label{thm:s1}
\begin{enumerate}
\item \tl{maynard\_S1\_expand\_W} (no hypotheses):
\[
\sum_{\substack{n<N\\ n\equiv v_0\,(W)}}\Big(\sum_{\substack{d\in D_k\\ d\mid n+h}}\lambda_d\Big)^2
=\sum_{d,e\in D_k}\lambda_d\lambda_e\,\#\{n<N:\ n\equiv v_0\ (W),\ d\mid n+h,\ e\mid n+h\}.
\]
\item \tl{tupleDvd\_lcm}, \tl{crt\_tuple\_count\_bounds}: the two divisibility
conditions are $[d_i,e_i]\mid n+h_i$; if the moduli $L_i$ are pairwise coprime
and coprime to $W$, the count is $N/(W\prod_iL_i)$ up to an error less than $1$.
\item \tl{count\_zero\_of\_not\_pairwise}, \tl{hpw\_of\_solvable}: if
$|h_i-h_j|\mid W$ for all $i<j$ and the entries are coprime to $W$, a pair $(d,e)$
whose moduli $[d_i,e_i]$ are not pairwise coprime contributes count $0$.
\item \tl{maynard\_S1\_error\_W}: under the hypotheses of (3), with the main term
over the remaining (``good'') pairs,
\[
\Big|S_1-\frac NW\sum_{(d,e)\ \mathrm{good}}\frac{\lambda_d\lambda_e}{\prod_i[d_i,e_i]}\Big|
\le\sum_{d,e\in D_k}|\lambda_d\lambda_e|.
\]
\end{enumerate}
\end{theorem}

\noindent Item (4) is the finite content of \cite[(5.2)]{May15}. Maynard's
restricted sum $\sum'$ is our sum over good pairs.

\begin{theorem}[\lbl{Lean-checked}; files \tl{MaynardInv.lean}, \tl{MuMult.lean}, \tl{MaynardDiag.lean}]\label{thm:diag}
Let $\lambda$ be supported on tuples in the box $\{1,\dots,R\}^k$ with
squarefree entries, and put
$y_u=\prod_i\mu(u_i)\varphi(u_i)\sum_{u\mid d}\lambda_d/\prod_id_i$, as in
\cite[Lemma~5.1]{May15}.
\begin{enumerate}
\item \tl{mobInv\_upSum}, \tl{upSum\_mobInv}: M\"obius inversion over
$k$-tuples in the box, both directions, for every $R$ and $k$.
\item \tl{maynard\_lam\_of\_y\_sqfree\_all} and
\tl{maynard\_y\_of\_lam\_sqfree\_all}: the formula
\[\lambda_d=\prod_i\mu(d_i)d_i\sum_{d\mid r}\frac{y_r}{\prod_i\varphi(r_i)}\]
and its inverse hold for every $R$ and $k$. The proof needs $\mu(a)\mu(x)=\mu(x/a)$ for
$a\mid x\mid c$, $c$ squarefree (\tl{muSign\_of\_sqfree}), which we derive from
the recursive definition of $\mu$.
\item \tl{maynard\_diag\_lcm}, \tl{maynard\_diag\_y\_sqfree}:
$\sum_{d,e}\lambda_d\lambda_e/\prod_i[d_i,e_i]=\sum_u y_u^2/\prod_i\varphi(u_i)$,
via Gauss's identity $\gcd(a,b)=\sum_{u\mid a,\,u\mid b}\varphi(u)$ in each
coordinate. No coprimality between coordinates is needed.
\item \tl{maynard\_S1\_diag\_W}: with the hypotheses of Theorem~\ref{thm:s1}(3),
\[
\Big|S_1-\frac NW\Big(\sum_{u}\frac{y_u^2}{\prod_i\varphi(u_i)}-B\Big)\Big|
\le\sum_{d,e}|\lambda_d\lambda_e|,\qquad
B=\sum_{(d,e)\ \mathrm{bad}}\frac{\lambda_d\lambda_e}{\prod_i[d_i,e_i]}.
\]
\end{enumerate}
\end{theorem}

\begin{remark}[two points the formalisation made explicit]
\begin{itemize}
\item The hypothesis ``the $[d_i,e_i]$ are pairwise coprime'' cannot be assumed
for the whole support: for $h=(0,2)$, $W=2$, the pair $d=(3,1)$, $e=(1,3)$ has
moduli $(3,3)$ (\tl{ex\_not\_hpw}), and on the full box the pair $d=e=(2,2,\dots)$
fails it (\tl{box\_hpw\_fails}). What holds is Theorem~\ref{thm:s1}(3): such
pairs have count $0$. Their main-term contribution $B$ need not vanish
($B=2/9$ in the example \tl{ex\_B\_eq}); Maynard controls it analytically, and
that bound is not formalised.
\item The change of variables needs squarefree support
(\tl{maynard\_inversion\_needs\_sqfree} gives a failing non-squarefree case).
\end{itemize}
\end{remark}

\section{What stays a hypothesis}\label{sec:hyp}

The following are \emph{not} formalised. To use Theorem~\ref{thm:skeleton}(4) or
(5) one must supply them. Labels refer to their status in the literature.

\begin{hypothesis}[the size of the finite errors]\label{hyp:err}
For Maynard's weights, $\sum_{d,e}|\lambda_d\lambda_e|$ and $B$ are
$o\big(\varphi(W)^kN(\log R)^k/W^{k+1}\big)$
(\lbl{Literature}, \cite[proof of Lemma~5.1]{May15}).
\end{hypothesis}

\begin{hypothesis}[the main term of $S_1$]\label{hyp:s1}
With $\lambda$ defined from a smooth $F$ supported on
$\{x\in[0,1]^k:\sum x_i\le1\}$ and $R=N^{\theta/2-\delta}$,
$\frac NW\sum_u y_u^2/\prod\varphi(u_i)=(1+o(1))\varphi(W)^kN(\log R)^kI_k(F)/W^{k+1}$
(\lbl{Literature}, \cite[Prop.~4.1, Lemmas~5.1 and~6.2]{May15}).
\end{hypothesis}

\begin{hypothesis}[the finite expansion and main term of $S_2$]\label{hyp:s2}
The analogue of Theorems~\ref{thm:s1} and~\ref{thm:diag} for
$S_2^{(m)}=\sum_n w_n\mathbf 1[n+h_m\ \text{prime}]$, and the asymptotic
$S_2=(1+o(1))\varphi(W)^kN(\log R)^{k+1}\sum_mJ^{(m)}_k(F)/(W^{k+1}\log N)$,
which uses a level of distribution $\theta$ for the primes (for $\theta<1/2$,
the Bombieri--Vinogradov theorem) (\lbl{Literature}, \cite[Prop.~4.1, Lemmas~5.2 and~6.3]{May15}; the level enters through the error term $O(y_{\max}^2N/(\log N)^A)$ of Lemma~5.2).
The finite expansion is elementary but not yet written in Lean; it is the
next step of this programme.
\end{hypothesis}

\begin{hypothesis}[the variational inequality]\label{hyp:var}
For $k=50$ there is $F$ on the enlarged region of \cite{Pol14b} with ratio above
$4$, together with the modified asymptotics that region needs
(\lbl{Literature}, \cite{Pol14b}). Our main paper re-certifies this inequality
exactly (\lbl{Computed}); the certificate is not in Lean.
\end{hypothesis}

\noindent Granting Hypotheses~\ref{hyp:err}--\ref{hyp:var}, with real constants
replaced by rationals on the correct side, the hypothesis of
Theorem~\ref{thm:skeleton}(5) holds and gives $\tl{BoundedGaps}\ 246$. That
implication is the content of \cite{May15,Pol14b}; we have formalised only the
finite part.

\section{Distance to twin primes}

The file \tl{Tuples.lean} also checks \tl{twinPrimeConjecture\_of\_DHL\_2}:
$\tl{DHL}\ 2$ implies the twin prime conjecture. This is a triviality, and it
is the only link between this note and twin primes. The route of
Theorem~\ref{thm:skeleton} needs a ratio criterion that sieve methods of this
kind cannot meet for $k=2$, because of the parity barrier: Polymath8b show,
by a modification of Selberg's parity argument, that their $k=3$ result under
GEH is the best possible from purely sieve-theoretic considerations
(\lbl{Literature}, \cite[abstract]{Pol14b}); see also the project's
parity-barrier paper. Nothing in this note goes beyond what is
already known.

\section{Reproducing the checks}\label{sec:audit}

\begin{itemize}
\item Toolchain: Lean~4.35.0 (official release), no Mathlib, no other dependencies.
\item \tl{cd lean \&\& lake build} builds all $21$ files of \tl{TwinPrimes/}
($5720$ lines; \tl{Reduction.lean}, the six \tl{Maynard*} files and the two \tl{Mu*} files total $2543$).
\item \tl{lake env lean Audit.lean} prints the axioms of $208$ declarations: $116$
use \tl{propext}, \tl{Classical.choice}, \tl{Quot.sound}; $53$ use
\tl{propext}, \tl{Quot.sound}; $11$ use \tl{propext}; $28$ use none. No
\tl{sorryAx}, no other axiom (run of 26 November 2026).
\item \tl{grep} finds no \tl{sorry}, \tl{native\_decide} or \tl{axiom}
declaration in \tl{TwinPrimes/}. Finite facts (admissibility of the $50$-tuple,
sign tables for small $c$) use \tl{decide}, checked by the kernel.
\end{itemize}

\appendix
\section{History of this document}\label{app:history}
\begin{description}
\item[1.0] (26 Nov 2026) First version. The one-dimensional Selberg files
(\tl{Selberg.lean}, \tl{SelbergEq.lean}, \tl{SelbergOpt.lean}; exact
diagonalisation and optimum for small and general $D$) are out of scope and
listed here only.
\end{description}

\begin{thebibliography}{9}
\bibitem{Axiom26} R.~Aanegola, \dots, J.~Thorner, \dots (Axiom Math),
\emph{Lean formalization of bounded gaps between primes} (PrimeGapsLib), 2026.
\url{https://primegaps.axiommath.ai/} (read 26 November 2026).
\bibitem{May15} J.~Maynard, \emph{Small gaps between primes}, Ann. of Math.\
\textbf{181} (2015), 383--413; arXiv:1311.4600v3 (read; SHA-256 of the PDF
\texttt{dce1a7a0\dots d63cda9}).
\bibitem{Pol14b} D.~H.~J.~Polymath, \emph{Variants of the Selberg sieve, and
bounded intervals containing many primes}, Res.\ Math.\ Sci.\ \textbf{1}
(2014), Art.~12; doi:10.1186/s40687-014-0012-7.
\end{thebibliography}

\end{document}
